\documentclass[main.tex]{subfiles}
\begin{document}
\section{Introduction}
\subsection{Question}
Determine the the relation between place, velocity, acceleration, and time in a model with constant acceleration.

\subsection{Dimensional analysis}
We will analyze this model:
$$
s=f(t,v,a)
$$
We define the dimension symbols
$$
\tau=t
$$
$$
\lambda=v\cdot t
$$
And then we define the dimensions of the physical quantities.
$$
a=\dfrac{\lambda}{\tau^2} = \dfrac{v\cdot t}{t^2}
$$
$$
v=\dfrac{\lambda}{\tau} = \dfrac{v\cdot t}{t}
$$
$$
s=\lambda=s
$$
$$
t=\tau=t
$$

Then we divide our physical quantities by their respective dimension and reduce our result.

\begin{align*}   
\dfrac{s}{\lambda}&=f\left(\dfrac{t}{\tau},\dfrac{v}{\lambda\tau^{-1}},\dfrac{a}{\lambda\tau^{-2}}\right)\\
\dfrac{s}{v\cdot t}&=f\left(\dfrac{t}{t},\dfrac{v}{vt\cdot t^{-1}},\dfrac{a}{vt\cdot t^{-2}}\right)\\
\dfrac{s}{v\cdot t}&=f\left(1,1,\dfrac{a}{vt^{-1}}\right)\\
\dfrac{s}{v\cdot t}&=f\left(\dfrac{at}{v}\right)
\end{align*}

So we end up with two $\Pi$-groups
$$
\Pi_1=\dfrac{s}{vt}\quad
\Pi_2=\dfrac{at}{v}
$$

and we will find the relation $f$ that connects them.





\subsection{Hypothesis}
We hypothesize that the final formula for the distance will be
$\mathbf{s(t) = \dfrac{1}{2} \cdot a \cdot t^2 + v_0 \cdot t}$
\end{document}
